\section*{Your turn}

Check out tag \code{ch11} of the companion repository and run \code{npm ci}
once. Exercises 1 to 4 need the database, both subgraphs and the router.
Exercises 5 to 8 need node, Docker and the .NET SDK, and nothing started by
hand: the cases script builds the sample, starts both of its subgraphs on ports
of its own, and stops them again.

\begin{minted}[fontsize=\footnotesize,breaklines]{text}
docker compose up -d mosaic-db
dotnet run --project src/Mosaic.Catalog     # 5101
dotnet run --project src/Mosaic.Api         # 5100
docker compose up -d mosaic-router          # 3002
\end{minted}

The two \code{dotnet run} commands want a terminal each. If port 3002 is
already taken, \code{docker compose --profile wire down} clears
chapter~\ref{ch:how-a-federated-query-actually-runs}'s router.

\begin{enumerate}
  \item \textbf{Watch the router fetch something nobody asked for.} Send this
  to port 3002 with the header \code{X-WG-Trace: true}:

  \begin{minted}[fontsize=\footnotesize,breaklines]{graphql}
{ browseProducts(first: 2) { nodes { title shippingCost { amount } } } }
  \end{minted}

  Read \code{extensions.trace.fetches} in the response. Find the document the
  router sent Catalog and the representations it sent Mosaic. Then take
  \code{shippingCost} out of your query and send it again. One word disappears
  from Catalog's document and the representations lose a field.

  \item \textbf{Make the subgraph fail the way the router never will.} Take one
  of those representations, delete its \code{category}, and post an
  \code{\_entities} query straight to port 5100 with three representations of
  which that is the middle one. Note which entity is null, which are not, and
  where the error's \code{path} points. Then look in Mosaic's console for the
  sentence the client did not get.

  \item \textbf{Break the setter.} Find \code{Category} on the \code{Product}
  stub in \code{Mosaic.Api}, change its \code{\{ get; set; \}} to
  \code{\{ get; \}}, and give it an initialiser so that it compiles. Restart
  Mosaic and ask the router for a shipping cost. Before you read the error,
  write down which of the two services you would have investigated first. Put
  the property back with \code{git checkout}.

  \item \textbf{Find out what a category is worth.} Ask the router for the
  storefront query at \code{first: 25} with \code{shippingCost} and again
  without it, three times each to warm both documents, reading Mosaic's
  timeline line after each. The resolver count moves by twenty-five and the SQL
  count does not move at all. Then work out what would have to change in
  \code{ProductPricingNode} for the second number to move as well.

  \item \textbf{Feel the difference a DataLoader makes.} Run

  \begin{minted}[fontsize=\footnotesize,breaklines]{text}
node scripts/entity-cases.mjs --print batch-is-one-lookup
node scripts/entity-cases.mjs --print naive-is-four-lookups
  \end{minted}

  and put the two logs side by side. The second one has four lookups, which is
  the obvious difference. Find the other one, and then find the line in
  \code{samples/entity-resolution} that causes it. If the script complains that
  a port is taken, an earlier run is still up.

  \item \textbf{Break a schema three ways.} Print each of these cases in turn:

  \begin{minted}[fontsize=\footnotesize,breaklines]{text}
external-unused    requires-not-external    external-nullability-wins
  \end{minted}

  Two of them print a composition error and one prints four lines and no error
  at all. Decide which of the three you would be most likely to ship, then read
  \code{scripts/entity-cases.mjs} to see what each one edited.

  \item \textbf{Give \code{@provides} something to lie about.} Print the case
  \code{provides-serves-what-it-stored} and read its four rows. Then open
  \code{Crates.cs} in the sample and change crate three's \code{PackedAs} so
  that it disagrees with Catalog too. Rerun the case. It fails, and the failure
  message is the assertion telling you that the sample no longer contains
  exactly one deliberate lie. No re-export is needed, because \code{PackedAs}
  is not in the schema. \code{git checkout} puts it back.

  \item \textbf{Decide a nullability, then test it.} Print the case
  \code{dangling-key-nulls-the-answer}. Then, in the same file, make
  \code{widgetByKey} return a nullable widget, and re-export that subgraph's
  schema over its file under \code{schema/samples/} before you rerun. Skip the
  re-export and the script stops before any case runs, because it checks each
  subgraph against its committed schema first. With the schema updated, the
  case fails on one assertion: the response no longer collapses. Read the new
  answer and decide which of the two you would rather hand a client.
  \code{git checkout} puts both files back.
\end{enumerate}

Exercise 3 is the one to do properly, and the reason is the sentence you write
down before reading the error. A missing setter compiles, passes composition,
and then breaks a field whose symptom points at the router and at the other
service. Where you look first is the whole cost of that bug, and writing the
guess down is the only way to find out afterwards whether you were right.
Chapter~\ref{ch:strangling-the-monolith} takes the remaining five domains out
of \code{Mosaic.Api}, and the arrangement it produces is four more copies of
what section~\ref{sec:ch11-null-at-the-boundary} just measured.
